\section{结论}
本文围绕“板凳龙”运动过程建立了统一的几何建模框架。首先，利用等距螺线弧长参数化与逐节递推，完成了盘入阶段全把手状态的精确恢复；随后，引入有限宽度矩形碰撞模型，求得无碰撞终止时刻；进一步以螺线与调头圆的交点为入点，在调头空间内构造两段相切圆弧；最后通过速度比例缩放，得到满足全队速度约束的最大恒定龙头速度。

本文的主要结论可以概括为四点：其一，等距螺线加弧长递推能够稳定描述板凳龙的连续运动；其二，有限宽度矩形模型比中心线模型更适合判定盘入阶段的碰撞边界；其三，Q3在角点到板凳轴线距离判据下给出的最小螺距约为$0.45$ m；其四，固定交点调头路径下的龙头最大恒定速度为$1.406$ m/s。

因此，本文给出的不是一组孤立数值，而是一套可复用的建模流程：先定路径，再定约束，最后在统一几何框架下完成搜索、判定与反推。这个流程与题目要求是一一对应的，也为后续类似板式编队运动问题提供了可迁移的思路。
